% chktex-file 1
% chktex-file 44
% chktex-file 45
% chktex-file 36
% chktex-file 3
% chktex-file 8

\documentclass[twoside]{article}
\usepackage{graphicx,fancyhdr,amsmath,amssymb,amsthm,subfig,url,hyperref}
\usepackage[margin=1in]{geometry}
\newtheorem{theorem}{Theorem}


%---------------- IMPORTANT: Student and Homework Information -------------

%%% PLEASE FILL THIS OUT WITH YOUR INFORMATION
\newcommand{\myname}{Muhammad Mujtaba}
\newcommand{\myid}{28100480}
\newcommand{\hwNo}{Assignment 3}
%%% END



\fancypagestyle{plain}{}
\pagestyle{fancy}
\fancyhf{}
\fancyhead[RO,LE]{\sffamily\bfseries\large LUMS}
\fancyhead[LO,RE]{\sffamily\bfseries\large CS 210 Discrete Mathematics}
\fancyfoot[LO,RE]{\sffamily\bfseries\large \myname: \myid @lums.edu.pk}
\fancyfoot[RO,LE]{\sffamily\bfseries\thepage}
\renewcommand{\headrulewidth}{1pt}
\renewcommand{\footrulewidth}{1pt}

%--------------------- This is the title of the document. DO NOT CHANGE IT ------------------------

\title{CS 210 \hwNo}
\author{\myname \qquad Student ID:\myid}
\newcommand{\currentdate}{3 September 2025}
\date{\currentdate}

%--------------------------------- AFTER Entering the Student and Homework Information, write your answers below  ----------------------------------

\begin{document}
\maketitle

\section*{Question 1}
Collaborated with: X, Y, Z.
\begin{enumerate}
    \item[(a)]  $3^5 = 243$
    \item[(b)]  No, a function f: A $ \rightarrow$ B cannot be injective because the cardinaltiy of A is greater than the cardinaltiy of B.
                In other words, there aren't enough elements in B to give each element of A a unique image.
    \item[(c)]  A = $\{1,2,3,4,5\}, B = \{a, b, c\}$ \\
                \quad f(1) = a, 
                \quad f(2) = a,
                \quad f(3) = b,
                \quad f(4) = b,
                \quad f(5) = c, \\
                Each element of B is hit by some value of A. Thus, surjective but not injective. 
\end{enumerate}

\section*{Question 2}
Collaborated with: X, Y, Z.
\begin{enumerate}
    
    \item[A.]  No, cos(x) is periodic \\ 
               Range = $[-1, 1]$
    \item[B.]  No, many integers share the same value \\
               All even x map to 1, all odd x map to -1 \\
               Range =$ \{1, -1\} $
    \item[C.]  No, for example f(2) is same as f(1/2) \\
               Deravative of f(x) = $ \frac{1 -x^2}{ (1+x^2)^2}$ and critical points are $\pm 1$ \\
               Thus, range = $[-\tfrac{1}{2}, \tfrac{1}{2}]$
    \item[D.]  No, every interval K that maps to the same integer k or k+ 1\\
               Range = Z All Integers
    \item[E.]  No, the function is even,  f(x) = f(-x) \\
               Range = $[0, \infty)$
\end{enumerate}

\newpage

\section*{Question 3}
Collaborated with: X, Y, Z.
\begin{enumerate}
    \item[(i)] 
    Let $f: Y \rightarrow Z \quad \text{and} \quad g: X \rightarrow Y $ \\
    Then their inverses are \\
    $f^{-1}: Z \rightarrow Y \quad \text{and} \quad g^{-1}: Y \rightarrow X $ \\
    Since,
    \[f(X) : \text{Y }\rightarrow \text{Z and g: X } \rightarrow \text{Y, we have f(g(X)) = Z}\]
    For the inverse, 
    \[(f \circ g)^{-1}(Z) = X\] 
    If we apply 
    \[f^{-1}(Z) = Y \quad \text{and then} \quad g^{-1}(Y) = X\] \\
    we get 
    \[X = g^{-1}(f^{-1}(Z))\]
    Hence, 
    \[(f \circ g)^{-1}(Z) = g^{-1}(f^{-1}(Z)) \]

    \item[(ii)] 
    The inverses are:
    \[f^{-1} = \{(a,1), (b,2), (c,3), (d,4)\},\]
    \[g^{-1} = \{(3,a), (1,b), (4,c), (2,d)\}.\]
\end{enumerate}

\section*{Question 4}
Collaborated with: X, Y, Z.
\begin{enumerate}
    \item[Case 1: ] If f is injective, then f is surjective \\
                    - Since f is injective, no 2 elements of A are mapped to the same element of B \\
                    - All n elements of A map to n unique elements in B \\
                    - But B itself has exactly n elements \\
                    - Therefore, every element of B must be used, nothing is left OUT \\
                    - Hence, f is surjective
    
    \item[Case 2: ] If f is surjective, then f is injective \\
                    - Since f is surjective, every element of B has atleast 1 element of A that maps into it \\
                    - Thus, the range(image) of f have n elements \\
                    - But A also has n elements in total \\
                    - So, each element of A map exactly one with each element of B \\
                    - Therefore, f must be injective
\end{enumerate}
So, for finite sets A and B of the same size, function is injective iff it is surjective

\newpage

\section*{Question 5}
Collaborated with: X, Y, Z.
\begin{enumerate}
    \item[(a)]  With n=4, that means each instruction has 4 - 1 = 3 choices. So, \\
                Total possible assignments = $3^{4}$ = 81  \\
                K is number of lower-index positions; S is the set of lower positions \\
                k = 3; S = (2,3,4) \\
                    - i = 1 (not lower): 3 choices \\
                    - i = 2 (lower): 1 choice \\
                    - i = 3 (lower): 2 choices \\
                    - i = 4 (lower): 3 choices \\
                    Total = 3 * 1 * 2 * 3 = 18 choices 

                Atleast 2 lower positions = Total choices - 3rd position
                \[81 -18 = 63\]

    
    \item[(b)]
        \begin{enumerate}
            \item[(i)]   a function? \\
                         Each assignment f is a function: \\
                         it maps every instruction to exactly one register. \\
                         Each element of the set of assignments is a set of function. \\
                         So, \\
                         Each valid Assignment is a function. The whole collection is a set of functions
            \item[(ii)]  injective? \\
                         Nope, Multiple instructions can target the same register, as long as the conditions are satisfied \\
                         Counter Example: \\
                            n = 4: assignment (2,1,1,1) is allowed and isn't injective
            \item[(iii)] surjective?? \\
                         Nope, An assignment cannot use all register \\
                         Same Counter Example: \\
                            n = 4: assignment (2,1,1,1) is allowed and isn' surjective as R3 and R4 isnt used
        \end{enumerate} 

    \item[(c)]  Yes, sometimes, but not always \\
                An injective function can exist when the assignment is a dearrangement and m $\le$ n -1; \\
                dearrangements exists for n $\ge$ 2 \\ \\
                Dearrangement is the first constraint... permutation with no fixed points \\
                For example:
                \[n = 4, m = 2\] 
                the permutation (4,1,2,3) or (2,1,4,3) are injective derangements and do satisfy the “at least 2 lower indices” requirement.

\end{enumerate}


\end{document}
